
\section{Modelos de comunicación}

\begin{frame}{Comunicación síncrona y asíncrona}
 \begin{itemize}
  \item Los procesos intercambian información
  para coordinar tareas en un sistema distribuido
  \cite{kazanavicius2023}.
  \item La comunicación puede ser
  \textbf{síncrona} o \textbf{asíncrona}
  \cite{kazanavicius2023}.
\end{itemize}

  {\tiny Fuente: \cite{kazanavicius2023}}
\end{frame}

\begin{frame}{Comunicación síncrona}
  \begin{itemize}
    \item El emisor ha realizado una solicitud y ahora está a la espera de una respuesta \cite{kazanavicius2023}.
    \item El proceso dependiente no continúa hasta que reciba el resultado.
    \item Podría provocar demoras si el receptor tarda mucho o no responde.
  \end{itemize}

  \vspace{0.3cm}
  \begin{block}{Ejemplo}
    Una tienda consulta si un pago fue aprobado y espera el resultado antes de confirmar la compra.
  \end{block}

  {\tiny Fuente: \cite{kazanavicius2023}}
\end{frame}

\begin{frame}{Comunicación asíncrona}
  \begin{itemize}
    \item El emisor envía un mensaje sin esperar el resultado de inmediato \cite{maharjan2023}.
    \item El receptor puede procesar el mensaje posteriormente.
    \item Este reduce la dependencia temporal entre los componentes \cite{kazanavicius2023}.
    \item Las colas permiten gestionar mensajes pendientes de procesamiento \cite{derlatka2023}.
  \end{itemize}

  {\tiny Fuentes: \cite{maharjan2023,derlatka2023,kazanavicius2023}}
\end{frame}

\begin{frame}{Comunicación síncrona frente a asíncrona}
  \small
  \begin{table}
    \centering
    \begin{tabularx}{\textwidth}{lXX}
      \toprule
      \textbf{Aspecto} & \textbf{Síncrona} & \textbf{Asíncrona} \\
      \midrule
      Respuesta & Espera el resultado & No espera el resultado inmediato \\
      Dependencia & Mayor dependencia temporal & Menor dependencia temporal \\
      Ventaja & Permite decidir con la respuesta & Permite continuar otras tareas \\
      Limitación & Puede generar esperas & Requiere gestionar mensajes \\
      Ejemplo & Consulta de pago & Cola de pedidos \\
      \bottomrule
    \end{tabularx}
  \end{table}

  \vspace{0.3cm}
  {\tiny Fuentes: \cite{kazanavicius2023,maharjan2023,derlatka2023}}
\end{frame}

\begin{frame}{Hilos, concurrencia y paralelismo}
  \begin{itemize}
    \item \textbf{Hilo:} es la unidad de ejecución dentro de un proceso.
    \item \textbf{Concurrencia:} dos o más hilos progresan durante un mismo período \cite{reinders2023}.
    \item \textbf{Paralelismo:} dos o más hilos se ejecutan simultáneamente \cite{reinders2023}.
  \end{itemize}

  \vspace{0.3cm}
  \begin{block}{Ejemplo}
  Tres trabajadores procesan pedidos diferentes.
  Si alternan su ejecución, existe concurrencia;
  si trabajan simultáneamente, existe paralelismo.
  \end{block}

  {\tiny Fuente: \cite{reinders2023}, cap. 1, pp. 28--29.}
\end{frame}


\begin{frame}{Modelos de comunicación}
  \begin{itemize}
    \item \textbf{Cliente-servidor:} un cliente solicita un servicio y otro proceso lo atiende \cite{zhang2023}.
    \item \textbf{Mensajería:} los procesos intercambian mensajes directamente o mediante intermediarios \cite{maharjan2023}.
    \item \textbf{RPC:} un proceso solicita la ejecución de una operación en otro proceso \cite{zhang2023}.
  \end{itemize}

  \vspace{0.3cm}
  \begin{block}{Aplicación en RabbitMQ}
    RabbitMQ actúa como intermediario entre productores y consumidores. También admite patrones de petición-respuesta \cite{maharjan2023}.
  \end{block}

  {\tiny Fuentes: \cite{zhang2023,maharjan2023}}
\end{frame}

\begin{frame}{Procesamiento asíncrono con RabbitMQ}
  \begin{enumerate}
    \item Un cliente realiza un pedido.
    \item La aplicación publica un mensaje en RabbitMQ.
    \item El broker dirige el mensaje a una cola.
    \item Un consumidor recibe y procesa el pedido.
    \item Otros consumidores pueden procesar pedidos diferentes.
  \end{enumerate}

  \vspace{0.3cm}
  \begin{block}{Ventaja}
    La aplicación puede continuar atendiendo solicitudes mientras los trabajadores procesan los pedidos.
  \end{block}

  {\tiny Fuentes: \cite{maharjan2023,kazanavicius2023,derlatka2023}}
\end{frame}

\begin{frame}{Aplicación al PFC: SREI}
\footnotesize

\begin{itemize}
    \setlength{\itemsep}{0.08cm}

    \item \textbf{Sistema:}
    Registro de Eventos Interculturales (SREI).

    \item \textbf{Servicio proveedor:}
    \texttt{servicio-eventos} (puerto 8080).
    Proporciona información de un evento intercultural.

    \item \textbf{Servicio consumidor:}
    \texttt{servicio-consulta} (puerto 8081).
    Solicita la información al proveedor y utiliza
    los datos recibidos para responder a la consulta.
\end{itemize}

\vspace{0.1cm}

\begin{block}{Implementación REST en Spring Boot}
\textbf{Proveedor (EventoController):}

\texttt{\detokenize{@GetMapping("/consulta")}}

\vspace{0.08cm}

\textbf{Consumidor (ConsultaController):}

\texttt{\detokenize{cliente.get().uri("/api/eventos/consulta")}}

\texttt{\detokenize{.retrieve().body(Map.class);}}
\end{block}

\vspace{0.1cm}

\begin{block}{Relación con la teoría}
El consumidor solicita información mediante HTTP
y espera la respuesta del proveedor para utilizarla.
Esto representa comunicación síncrona mediante REST.
\end{block}

\end{frame}